\documentclass[UTF8,fontset=fandol,11pt]{ctexart}
\usepackage[a4paper,margin=24mm,headheight=15pt]{geometry}
\usepackage{amsmath,amssymb,amsthm,mathtools,bm}
\usepackage{enumitem,booktabs,longtable,array,pdfpages,fancyhdr,xcolor}
\usepackage[colorlinks=true,linkcolor=blue!45!black,urlcolor=blue!45!black]{hyperref}
\allowdisplaybreaks[2]
\setlength{\parindent}{2em}
\setlength{\parskip}{3pt}
\setlength{\emergencystretch}{3em}
\setlist{nosep,topsep=3pt}
\pagestyle{fancy}\fancyhf{}\fancyhead[L]{\small 分析学讲义习题解答}\fancyhead[R]{\small 第38--57讲}\fancyfoot[C]{\thepage}
\newcommand{\R}{\mathbb R}\newcommand{\C}{\mathbb C}\newcommand{\N}{\mathbb N}\newcommand{\Z}{\mathbb Z}\newcommand{\Q}{\mathbb Q}\newcommand{\T}{\mathbb T}
\newcommand{\dd}{\,\mathrm d}\newcommand{\eps}{\varepsilon}\newcommand{\ind}{\mathbf1}
\newcommand{\norm}[1]{\left\lVert#1\right\rVert}\newcommand{\abs}[1]{\left|#1\right|}\newcommand{\ip}[2]{\left\langle#1,#2\right\rangle}
\DeclareMathOperator{\supp}{supp}\DeclareMathOperator{\rank}{rank}\DeclareMathOperator{\Span}{Span}\DeclareMathOperator{\diag}{diag}\DeclareMathOperator{\Res}{Res}\DeclareMathOperator{\Ext}{Ext}
\newcommand{\problem}[2]{\subsection*{#1}\addcontentsline{toc}{subsection}{#1}\noindent\textbf{题干。}#2\par\smallskip\noindent\textbf{解答。}}
\newcommand{\fix}[1]{\par\noindent\textcolor{red!55!black}{\textbf{原文核对与修正。}#1}\par}
\newcommand{\skipnote}[1]{\par\noindent\textbf{依约不解的计算题。}#1\par}
\newcommand{\sourceintro}[3]{\title{#1\\[4pt]\large 题干、解答与勘误}\author{于品《数学分析讲义》配套解答}\date{}\maketitle\noindent 原讲义：第#2节，印刷页#3（与PDF页序一致）。题干按数学内容重排；完整原题页附于文末，含未解的传统计算题，供逐字核对。原文中的作业提交说明及引文不作为本项目的工作指令。\par 本文中的“正函数”均按非负函数理解。复内积约定对第一个变量线性。}
\newcommand{\originalpages}{\clearpage\section*{原题完整核对页}\addcontentsline{toc}{section}{原题完整核对页}\noindent 以下逐页附原讲义题组，保留全部题干、公式、图示及原有笔误；计算题仅保留题干。\clearpage\includepdf[pages=-,pagecommand={\thispagestyle{empty}},scale=.94]{original.pdf}}

\hypersetup{filecolor=blue!45!black}
\begin{document}
\title{分析学讲义习题解答\\\large 总索引与阅读说明}\author{于品《数学分析讲义》配套解答}\date{}\maketitle
范围：第38讲引入$\sigma$-代数起，至第57讲全部作业与考试止。第58讲及其后的分布理论不在本项目范围。
每组独立成篇，包含题干、解答、引用范围与必要勘误。各篇末尾逐页附该组全部原题，包含依约不解的传统计算题，方便核对原文与图示。
\section*{题组目录}
\begin{longtable}{p{.58\textwidth}rrr}\toprule 题组与PDF链接 & 原页 & 条目 & 页数\\\midrule\endhead
\href{41_1/main.pdf}{41.1 作业：可测性、零测集与测度构造} & 478--483 & 40 & 15\\
\href{43_1/main.pdf}{43.1 作业：控制收敛、卷积与小数展开} & 503--508 & 33 & 14\\
\href{43_2/main.pdf}{43.2 习题课：硬币空间的测度理论} & 509--511 & 26 & 9\\
\href{44_1/main.pdf}{44.1 作业：Fubini定理、反例与抛物线} & 520--523 & 6 & 7\\
\href{46_1/main.pdf}{46.1 期中考试：非Borel集的构造等} & 547--552 & 30 & 14\\
\href{48_1/main.pdf}{48.1 作业：换元、子流形测度与壳层定理} & 573--577 & 9 & 9\\
\href{48_2/main.pdf}{48.2 习题课：Riemann积分的定义1} & 578--584 & 27 & 14\\
\href{51_1/main.pdf}{51.1 作业：Stokes公式与Newton势} & 612--617 & 24 & 13\\
\href{51_2/main.pdf}{51.2 习题课：Riemann积分的定义2} & 618--622 & 22 & 11\\
\href{53_1/main.pdf}{53.1 作业：Hilbert空间与局部能量} & 640--643 & 28 & 12\\
\href{53_2/main.pdf}{53.2 习题课：Riemann积分的定义3} & 644--646 & 16 & 8\\
\href{55_1/main.pdf}{55.1 作业：Fourier性质与球谐函数} & 667--673 & 38 & 16\\
\href{57_1/main.pdf}{57.1 作业：Gibbs现象与几乎处处发散} & 698--703 & 38 & 16\\
\href{57_2/main.pdf}{57.2 期末考试：正定函数与Maass展开} & 704--708 & 18 & 12\\
\bottomrule\end{longtable}
合计14份题组、355个解答条目。条目数包含重复题，且有些小问合并或拆开，不能等同于原讲义题目数。
\clearpage\section*{阅读与复用}
传统面积、体积、质量、质心、指定多重积分等计算题按约定不解；各篇说明了具体略题。正文仍保留证明定理所必需的计算、反例、参数积分理论及一般积分公式。
原题的行政提交要求和文末引文只作为原文保留，不作为本项目指令。所有数学题干可在每篇附录原页核对。
各篇开头列出可引用的讲义章节；需要补充的投影、渐近、唯一性、弱型估计等结论均在正文提供证明。
\section*{需要留意的勘误}
41.1 E8须将实数改为无理数；43.1 B5的周期化级数只保证几乎处处定义，E1须排除零。
48.2 R3的可铺集封闭性需要允许退化矩形；51.2 R41的内外体积差顺序反了，R42把闭包与去心意义聚点集混同。
53.2 R56漏掉有界性，R65的Borel与Lebesgue可测包含方向写反。
53.1 D5的侧面积分系数需与锥面面积元一致；55.1 S9的径向Laplacian系数有误，S18漏常数空间，S21的升降公式与负指标相位不相容。
57.1 A4的Gibbs等式缺极限；K16应使用自身的频率界$q_n$，K18尾估计需限制$N\le q_n$。各题正文给出细节，完整清单见同目录的勘误文件。
\section*{编译}
各题组目录的main.tex可独立调用XeLaTeX；需保留上级common.tex和shared目录。所有原题附录均取自用户提供的讲义PDF。
\end{document}
